\section{Finite deadlines and post-repair information}
\label{sec:finite-deadline}

\subsection{An explicit finite protocol}
Fix integers $N,R,H$ with $N+R<H$. The asymptotic policy uses
$N=R=\lfloor H_\delta/3\rfloor$, but a finite experiment must specify its
budgets explicitly. Run the polynomial race of
Theorem~\ref{thm:polynomial} as follows.
\begin{quote}\ttfamily\small
For n = 0, ..., N:\\
\hspace*{1em}Check decoder eligibility using the first n samples.\\
\hspace*{1em}If eligible, select the first decoder and leave this loop.\\
\hspace*{1em}If n = N, charge the old-state tail and declare a fallback label.\\
\hspace*{1em}Otherwise collect one more pre-repair observation.\\
After selecting a decoder:\\
\hspace*{1em}Attempt repair at most R times, stopping on success.\\
\hspace*{1em}On success, classify the post class and apply the decoder.\\
\hspace*{1em}On all failures, charge the old-state tail and declare fallback.
\end{quote}
The boundary sample $n=N$ is checked before timeout. In the separated
version the selected decoder is constant and its label is retained on every
subsequent branch, including failed repair.

Let $T$ be the untruncated race time and $f=(1-q)^R$. The same argument as
in Section~\ref{sec:proofs} gives the finite error bound
\begin{equation}
\PP_i(\widehat\theta\ne\theta_i)
\le \frac{\delta}{4}+\PP_i(T>N)+f+(M-1)\rho^{H-N-R},
\label{eq:finite-error-general}
\end{equation}
when maximum-likelihood post classification is used. The last term is
omitted for one class. It also gives
\begin{equation}
\EE_iC\le
g_i\EE_i(T\wedge N)+\frac{\kappa_i(1-f)}q
+g_iH\{\PP_i(T>N)+f\}.
\label{eq:finite-cost-general}
\end{equation}
These are unconditional bounds. In particular, a path-specific sample count
must not be mixed with an unconditional timeout probability and called a
conditional expectation. A positive threshold parameter alone does not imply
that arbitrary finite budgets achieve error $\delta$.

For Bernoulli post means, the implementation may instead use nearest-mean
classification. With minimum distinct-mean gap $\Delta_Q$, its last error
term can be replaced by
$2\exp(-(H-N-R)\Delta_Q^2/2)$, by the bounded-sum inequality
\citep{hoeffding1963}. Both classifiers coincide for the symmetric two-class
means $0.8$ and $0.2$ used in the main finite experiment.

\subsection{A quantitative information bound before taking a limit}
The exact equality partition describes what can eventually be distinguished.
At a finite deadline, the divergences between post laws also matter.
Let $N_P,N_Q$ be the numbers of acquired pre- and post-repair observations,
respectively, and let
\[
h_\delta=\operatorname{kl}(1-\delta,\delta)
=(1-2\delta)\log\frac{1-\delta}{\delta},
\qquad 0<\delta<1/2.
\]

\begin{theorem}[Finite-deadline information lower bound]
\label{thm:deadline-lower}
For every uniformly $\delta$-correct policy in the ordered experiment with
deadline $H$, and every opposite-old-label pair $i,j$,
\begin{equation}
\EE_iN_P\,\KL(P_i\Vert P_j)
+\EE_iN_Q\,\KL(Q_i\Vert Q_j)\ge h_\delta.
\label{eq:two-phase-info}
\end{equation}
In particular,
\begin{equation}
J_i^*(\delta;H)\ge
g_i\max_{j:\theta_j\ne\theta_i}
\frac{[h_\delta-H\KL(Q_i\Vert Q_j)]_+}
     {\KL(P_i\Vert P_j)}.
\label{eq:deadline-cost-lower}
\end{equation}
The latter is a lower bound, not a general matching coefficient.
\end{theorem}
\begin{proof}
Include the policy randomization in the complete transcript. It has the same
distribution under both environments. Conditional on an observed history,
the policy uses the same action kernel; each acquired pre observation
contributes $\KL(P_i\Vert P_j)$ to the expected log-likelihood ratio, and each
post observation contributes $\KL(Q_i\Vert Q_j)$. Repair-success trials,
deterministic time advancement, and environment-independent policy coins
contribute zero. The finite-horizon chain rule therefore gives the left side
of \eqref{eq:two-phase-info} as the transcript divergence. Ignoring an available
observation can only reduce what the terminal rule uses, and does not
invalidate the bound.

Map the transcript to the event that the terminal answer equals $\theta_i$.
Under $i$ this event has probability at least $1-\delta$; under the
opposite-label $j$ it has probability at most $\delta$. Data processing for
relative entropy and monotonicity of binary divergence in these two ranges
give the lower bound $h_\delta$. This is the standard sequential
change-of-measure principle \citep{kaufmann2016}, applied to the two phases.
Finally $N_Q\le H$, $C\ge g_iN_P$, and nonnegativity give
\eqref{eq:deadline-cost-lower}. Take the infimum over uniformly correct
policies.
\end{proof}
This bound deliberately uses opposite-label alternatives because it concerns
the \emph{final old answer}. It need not recover the sharper decoder bound at
a fixed post experiment. Their purposes differ: \eqref{eq:deadline-cost-lower}
quantifies a finite information budget, while Theorem~\ref{thm:main} exploits
which answer can be selected before a fixed post experiment.

\subsection{A shrinking-channel limit that removes the coefficient advantage}
The fixed-instance qualification in Theorem~\ref{thm:main} cannot be dropped.
Allow only the post laws to vary with $\delta$, writing $Q_i^\delta$, while
the pre laws, labels, costs, and $q>0$ remain fixed.

\begin{theorem}[Collapse to the separated coefficient]
\label{thm:collapse}
Let $H_\delta=\lceil L_\delta^2\rceil$. If, for every opposite-label pair,
\begin{equation}
H_\delta\KL(Q_i^\delta\Vert Q_j^\delta)=o(L_\delta),
\label{eq:shrinking-condition}
\end{equation}
then the optimal joint cost for this sequence of experiments satisfies
\begin{equation}
\lim_{\delta\downarrow0}\frac{J_i^*(\delta;Q^\delta)}{L_\delta}
=\frac{g_i}{D_i^{\rm sep}}\quad\text{for every }i.
\label{eq:collapse}
\end{equation}
Thus post laws can be pairwise distinct at every confidence level, yet give
no improvement in the leading coefficient along the sequence.
\end{theorem}
\begin{proof}
Since $h_\delta/L_\delta\to1$, \eqref{eq:deadline-cost-lower} and
\eqref{eq:shrinking-condition} imply
\[
\liminf_{\delta\downarrow0}\frac{J_i^*(\delta;Q^\delta)}{L_\delta}
\ge g_i\max_{j:\theta_j\ne\theta_i}\frac1{\KL(P_i\Vert P_j)}
=\frac{g_i}{D_i^{\rm sep}}.
\]
For the matching upper bound, run the separated pre-repair race from
Section~\ref{sec:proofs}. Its constant decoder commits to the old label;
post observations can be ignored. Its error analysis therefore depends only
on the fixed pre laws. Repair success and fully charged tails still have the
same fixed $q$, so its expected cost has coefficient $g_i/D_i^{\rm sep}$.
This policy is also joint-feasible for every $Q^\delta$. Combining bounds
proves the result.
\end{proof}

For a concrete example, retain the three pre means and labels of
\eqref{eq:example} and set, for $L_\delta\ge1$,
\begin{equation}
Q_i^\delta=\Bern\left(\frac12+\frac{a_i}{L_\delta}\right),
\qquad (a_1,a_2,a_3)=(-1/8,0,1/8).
\label{eq:shrinking-example}
\end{equation}
The post means are always distinct and lie in $[3/8,5/8]$.
The elementary inequality
$\KL(\Bern(u)\Vert\Bern(v))\le (u-v)^2/[v(1-v)]$
follows from $\log x\le x-1$. Hence all post divergences are
$O(L_\delta^{-2})$, and \eqref{eq:shrinking-condition} holds.

If any one of these post experiments is frozen and the confidence level is
then sent independently to zero, its singleton partition admits a universal
decoder and Theorem~\ref{thm:main} gives bounded expected cost. Along
\eqref{eq:shrinking-example}, however, \eqref{eq:collapse} gives the positive
separated coefficient. The two limiting operations therefore have different
outcomes. This is a quantified boundary of the original repair theorem, not
a contradiction to it.
